%------------------------------------------------------------------------------%
\documentclass[fleqn,utf8,aspectratio=169,12pt]{beamer}

\usepackage{gaal}

\title{Ângulo entre Planos}

%------------------------------------------------------------------------------%
\begin{document}

\addtitle

%------------------------------------------------------------------------------%
\section{Ângulo entre planos}

%------------------------------------------------------------------------------%
\begin{frame}{Ângulo entre dois planos}
  Se
  \(
  \vec{n}_{1}
  \)  e
  \(
  \vec{n}_{2}
  \)
  são vetores normais aos planos

  \pause
  O ângulo entre dois planos é
  \pause
  \[
    \cos{\theta}
    \pause
    =\frac{\abs{\vec{n}_{1}\cdot\vec{n}_{2}}}
    {
      \norm{\vec{n}_{1}}
      \norm{\vec{n}_{2}}
    }
  \]
  \pause
  O módulo permite escolher o menor ângulo entre os planos
\end{frame}

%------------------------------------------------------------------------------%
%------------------------------------------------------------------------------%
\begin{question}
Determine o ângulo entre os planos
  \[
    \alpha\colon x + y + z - 1 = 0
  \]
  \[
    \beta\colon 2x - y + 2z + 3 = 0
  \]
\end{question}

%------------------------------------------------------------------------------%
\begin{resolution}{Solução}
  Vetores normais
  \[
    \vec{n}_{1} = (1,  1, 1)
    \qquad\qquad
    \vec{n}_{2} = (2, -1, 2)
  \]

  \[
    \pause
    \vec{n}_{1}\cdot\vec{n}_{2}
    \pause
    = 2 - 1 + 2
    \pause
    = 3
  \]

  \[
    \pause
    \norm{\vec{n}_{1}}
    \pause
    = \sqrt{3}
  \]

  \[
    \pause
    \norm{\vec{n}_{2}}
    \pause
    = \sqrt{4 + 1 + 4}
    \pause
    = 3
  \]

  \[
    \pause
    \cos{\theta}
    \pause
    = \frac{3}{3\sqrt{3}}
    \pause
    = \frac{1}{\sqrt{3}}
  \]

  \[
    \pause
    \theta
    = \arccos{\left(\frac{1}{\sqrt{3}}\right)}
  \]
\end{resolution}

%------------------------------------------------------------------------------%
\endinput

%%------------------------------------------------------------------------------%
%\section{Lista Mínima}
%
%%------------------------------------------------------------------------------%
%\begin{frame}{Lista Mínima}
%  \begin{enumerate}
%    \item
%  \end{enumerate}
%
%  \vfill
%  Atenção: A prova é baseada no livro, não nas apresentações
%\end{frame}

%------------------------------------------------------------------------------%
\end{document}
%------------------------------------------------------------------------------%
